\clearpage
\chap{A reader's workshop}\label{app:workshop}
These short problems isolate choices that recur in the main development. Each has a complete solution. They are reading exercises and direct applications of the cited theory, not additional entries in the formal theorem inventory.

\subsection*{1. A stable limit does not make every seed persistent}
At every level take \(X_n=\{0,1\}\), with each bonding map sending both elements to zero. Find the inverse limit and the image of evaluation at every level. Does nonemptiness of every finite fiber make the seed one extendable?

\textbf{Solution.} For a coherent family, \(x_n=\rho_n(x_{n+1})=0\) for every \(n\). The all-zero family is coherent and is therefore the only point. Every evaluation image is \(\{0\}\). Equivalently, the image of restriction from any strictly later level is already \(\{0\}\), so the persistent image is \(\{0\}\). The seed one never extends. This is compatible with Proposition~\ref{KS58}, which proves nonemptiness and characterizes the persistent image, not surjectivity of every evaluation.

\subsection*{2. A residual chart can change without changing the defect}
Let \(d:\R^2\to\R\) be \(d(u,r)=r\). For \(a\in\R\), let \(A_a(r)=(ar,r)\). Compute the exact repair \(P_a\), compare it with \(P_b\), and decide when they commute.

\textbf{Solution.} Since \(dA_a=I\),
\[
 P_a(u,r)=(u-ar,0),\qquad
 P_b(u,r)-P_a(u,r)=((a-b)r,0).
\]
Both fix the horizontal exact axis. Thus \(P_aP_b=P_b\) and \(P_bP_a=P_a\). They commute exactly when they are equal, which occurs exactly when \(a=b\), as can be checked at input \((0,1)\). Every repair has zero residual. The different answers reflect the different right inverses, as in Theorem~\ref{KS65}.

\subsection*{3. Pairwise invertibility does not supply a triple law}
Let \(t\) be a nonidentity transposition in \(S_3\). Suppose three proposed overlap comparisons require frames \(b_0,b_1,b_2\) satisfying
\[
 b_1=tb_0,\qquad b_2=tb_1,\qquad b_2=tb_0.
\]
Every comparison is individually invertible. Can the three frames exist?

\textbf{Solution.} The first two equations give \(b_2=t^2b_0=b_0\). The third then gives \(b_0=tb_0\), and right cancellation gives \(t=1\), contrary to the assumption. Invertibility of every proposed overlap map does not prove the composition law needed for descent. Theorem~\ref{KS34} applies to data satisfying that law, not to arbitrary lists of pairwise isomorphisms.

\subsection*{4. Zero noisy output is not a certificate of reference exactness}
Take \(D(a,b)=a\), \(\widetilde D(a,b)=a+\tfrac1{10}b\), and input \(x=(-10,100)\). The measured operator returns zero exactly. Find the distance to the reference kernel and evaluate the reference uncertainty bound.

\textbf{Solution.} The reference kernel is the vertical axis, so the distance is ten. Its residual gain is one, and the operator error norm is \(1/10\). With \(y=0\) and \(\eta=0\), Proposition~\ref{KS69} gives
\[
 \dist(x,\ker D)\le\frac1{10}\sqrt{10^2+100^2}=\sqrt{101}.
\]
The true distance ten is below this bound. A formula omitting the operator-error term would instead claim a zero distance, which is false. The input is exact for the measured model and not for the reference one.

\subsection*{5. A prefix computation can hide a wraparound defect}
On a four-cycle, take proposed differences \((2,-1,3,-3)\) and anchor zero. Integrate the first three edges. Which equation fails, and what changes when only the final edge is allowed to move?

\textbf{Solution.} The first three edges give values \((0,2,1,4)\). Their required last difference is \(x_0-x_3=-4\), not \(-3\). The total prescribed sum is one, so the original loop has no realization. Keeping the first three edges forces the unique correction \((2,-1,3,-4)\), whose sum is zero. This is a protected repair of the edge data. It is not a realization of the original data without modification.

\subsection*{6. Residual decay can hide unprotected drift}
In maximum-norm $\R^3$, use $p(s,u,r)=s$ and $d(s,u,r)=r$. Compare the paths $(1,0,(n+1)^{-1})$ and $(1,H_n,(n+1)^{-1})$, where $H_n=\sum_{j=1}^{n}1/j$. Which hypotheses distinguish them?

\textbf{Solution.} Both preserve $p=1$ and have the same residual ratios $(n+1)/(n+2)$. The first moves by $1/((n+1)(n+2))$ and converges to $(1,0,0)$. The second moves by $1/(n+1)$ and has an unbounded second coordinate. In the variable-rate theorem, the first has $C_n=1/(n+2)$ and a summable weighted budget. The second has $C_n=1$, making that budget harmonic and divergent. The difference is movement, not residual decay or semantic preservation.

\subsection*{7. How a finite observation changes the apparent rate}
Let $\eta_n=1/(n+1)$. Find its profile after sampling at $\phi(k)=2^k-1$. Does the new profile prove exponential convergence in the original index?

\textbf{Solution.} The sampled value is $\eta_{\phi(k)}=2^{-k}$. Recovering the original index gives $k=\log_2(n+1)$ on the sampled stages, so the value is still $1/(n+1)$. The schedule preserves vanishing but changes the meaning of the horizontal axis. No faster original process has been proved. Proposition~\ref{KS75} is a convergence equivalence under a monotone cofinal schedule, not a rate-invariance theorem.

\subsection*{8. What does a zero floor fail to certify?}
For the diagonal model of Section~\ref{sec:escape}, can an exact protected answer exist in $\ell^2\oplus\R$? Does a coherent family of exact finite prefixes exist? Explain why the answers differ without contradicting finite-fiber existence.

\textbf{Solution.} Exactness with protected value one forces every sequence coordinate to equal minus one. This is not square summable, so there is no protected Hilbert answer. Each finite prefix nevertheless has the unique solution $(-1,\ldots,-1;1)$, and these solutions are coherent under truncation. Their inverse limit contains the all-minus-one sequence in the full product. The finite-fiber theorem supplies that coherent point. It does not prove that the point is in the image of the square-summable semantic space. The correction image is dense but does not contain the raw residual, so its floor is zero and its algebraic class is nonzero.

\subsection*{9. A phase witness without enumerating every gauge}
For $c(x,y)=x_2y_1$ on $\mathbf F_2^2$, show that no scalar coboundary equals $c$. Then compute the square of the lifted label $(0,e_1+e_2)$ in the eight-element extension.

\textbf{Solution.} Scalar coboundaries $a(x)+a(y)-a(x+y)$ are symmetric, but $c(e_1,e_2)=0$ and $c(e_2,e_1)=1$. Thus no scalar coboundary equals $c$. For $x=e_1+e_2$, $c(x,x)=1$, so $(0,x)^2=(1,0)$ and its fourth power is the identity. The corresponding matrix is $XZ$, whose square is $-I$. The additional phase bit records precisely the sign lost by projecting back to the four-label group.

\subsection*{10. The finite energy identity}
In $\R^2$, first prescribe $a_1+a_2=1$, then also prescribe $a_1=3/4$. Find the two minimum norm corrections and check the exact nested-affine identity.

\textbf{Solution.} The first correction is $(1/2,1/2)$, with energy $1/2$. The second is $(3/4,1/4)$, with energy $5/8$. Their difference has squared norm $1/8=5/8-1/2$. The identity uses the affine nature of both constraint sets. For the nested closed convex sets $[1,\infty)$ and $[2,\infty)$ in $\R$, the minimum points instead have squared difference one and energy difference three.

\subsection*{11. A basis is not automatically an isometry}
In Euclidean $\R^2$, let $p(a)=a_1-2a_2$, $d(a)=a_1$, $r=1$, and $\alpha=1$. Compute the protected regularized correction using the basis $(2,1)$ of the kernel. What changes if the Gram matrix is omitted?

\textbf{Solution.} The coordinate $c$ gives $a=(2c,c)$, with $S=2$ and $G=5$. The normal equation is $9c=2$, so $a=(4/9,2/9)$. Its objective is $5/9$. Replacing $G$ by one gives $c=2/5$ and $a=(4/5,2/5)$, whose actual ambient objective is $21/25>5/9$. That replacement solves a different metric problem.

\subsection*{12. A robust negative certificate}
A nominal scalar system has $T=1$, $r=1$. The operator error is at most $1/10$ and the data error at most $1/20$. Can a correction of norm at most $1/2$ achieve true residual at most $3/10$? What does the same test say at budget two?

\textbf{Solution.} At budget $1/2$, the nominal optimum is $1/2$. Subtracting $1/20+(1/10)(1/2)$ leaves $2/5>3/10$, so the proposed accuracy is excluded for every admissible true model. At budget two, the lower estimate is negative and clips to zero. It excludes nothing. The disappearance of this certificate is not proof of exact feasibility.

\subsection*{13. A smaller coordinate residual without a better correction}
For the protected example in Section~\ref{sec:inexact-repair}, keep the same correction $(1/4,1/4)$ but replace the kernel basis $B=(1,1)^{\mathsf T}$ by $B'=B/100$. Compute the coordinate normal residual and its actual Hilbert norm.

\textbf{Solution.} In the original basis, $c=1/4$, $G=2$ and $q=3c-1=-1/4$. In the new basis, $c'=25$, $G'=1/5000$ and $q'=-1/400$. The coordinate residual is one hundred times smaller, but
\[
 (q')^{\mathsf T}(G')^{-1}q'=\frac{5000}{160000}=\frac1{32}
 =q^{\mathsf T}G^{-1}q.
\]
The physical correction and the certificate have not improved. The changed number records a changed basis.

\subsection*{14. How small must the solver residual be?}
For $n\ge1$, suppose $\alpha_n=n^{-2}$, the measured-data error is bounded by $n^{-2}$, and the exact-data regularized corrections converge to the desired minimum-norm answer. Does a solver normal-residual tolerance $n^{-2}$ guarantee convergence of the proposals? Give a sufficient replacement.

\textbf{Solution.} The noise term is at most $1/(2n)$, but the solver term is at most one. Thus the displayed tolerance does not provide the required guarantee. The scalar counterexample $T=0$, $r=0$, $a_n=1$ meets it and does not approach the exact optimizer. Replacing it by $n^{-3}$ makes the solver term at most $1/n$, so the triangle inequality and the assumed vanishing bias prove convergence. This is a sufficient worst-case rule, not a necessary condition for every particular operator.

\subsection*{15. Exact optimization of measured data is not exact realization}
Let $d$ be the identity on $\R$, protect no nonzero observable, take raw input $f=2$, measured residual $r=1$, and $\alpha=1$. Find the exact measured-data optimizer and compare the fit residual with the actual repaired residual.

\textbf{Solution.} The objective $(a-1)^2+a^2$ is minimized at $a=1/2$. Its normal residual is zero, so the solver certificate is exact. The measured fit residual is $r-da=1/2$. The actual repaired candidate is $f-a=3/2$, whose model residual is $3/2$. The difference is $df-r=1$. A certificate about solving the measured optimization problem does not remove measurement error or the regularization bias.

\subsection*{A checklist for applying the theory}
Before selecting a theorem, state the raw and exact objects, the coefficient structure, the maps that relate views, and the data that must remain fixed. Then ask which conclusion is needed: existence, a description of all answers, an identification-preserving comparison, a stable distance estimate, or an executable finite certificate. Finally record the hypotheses that make that conclusion available.

The examples above fail only when one of those distinctions is suppressed. The missing information is not supplied by giving the same construction a more general name.

